\chapter{Полное доказательство теоремы о сходимости AFD}\label{app:afd-proof}

В настоящем приложении приводится развёрнутое доказательство
теоремы~\ref{thm:afd} главы~\ref{ch:afd} о
\emph{безусловном} выполнении условия~\eqref{eq:cond14} для алгоритма
AFD в схеме VIJI-Restarted и о согласующейся с этим
локально-оптимальной скорости $\mathcal O(L_1 D^3 / T^{3/2})$. Контроль
направленной ошибки достигается одной дополнительной конечно-разностной
(FD) секущей на каждой внутренней итерации.

\section{Алгоритм AFD: формализация}
\label{sec:app-afd:alg}

Зафиксируем входные данные внутренней итерации сегмента $s$:
\begin{itemize}
\item текущая итерация $v_k$ и матричный приближённый якобиан $B_k$;
      обозначим $J_k := \nabla F(v_k)$ и определим
      $M_0 := \tilde M_{0} + L_1 D \sqrt{d}$, где $\tilde M_{0}$~---
      Frobenius-оценка $\Fro{B_k^{(r)} - \Jstar} \le \tilde M_{0}$
      из~\eqref{eq:sp-bounded}/\eqref{eq:L0-warm}, $L_1 D \sqrt{d}$
      покрывает $\Fro{\Jstar - J_k}\le L_1\Norm{v_k - x^*}\sqrt{d}\le L_1 D\sqrt{d}$
      (триангуляция), $D$~--- диаметр $X$. Тогда $\Fro{B_k^{(r)} - J_k}\le M_0$
      равномерно по $k, s, r$;
\item регуляризатор $\beta_k \ge \beta_{\min} > 0$ из правила VIJI
      (см.~Algorithm~1 в~\cite{Agafonov2024VIJI});
\item параметры FD: шаг $\tau \in (0, 1]$ (по умолчанию $\tau = 1/8$),
      порог близости $\gamma > 0$ (определяется ниже), максимальное
      число refinement-шагов $r_{\max} \in \mathbb{N}$~--- гиперпараметр
      алгоритма (трункация цикла).
\end{itemize}
AFD-итерация задаётся как фиксированная последовательность шагов
($r$~--- индекс refinement-шага):
\begin{algorithm}[t]
\caption{AFD: одна внутренняя итерация в $v_k$}
\label{alg:afd-inner}
\begin{algorithmic}[1]
\State Compute candidate step $d^{(0)}_k \;=\; -(B_k + \beta_k I)^{-1} F(v_k)$.
\For{$r = 0, 1, \ldots, r_{\max}$}
  \State Compute FD direction
         $g^{(r)}_k \;=\; \tfrac{1}{\tau}\bigl(F(v_k + \tau\, d^{(r)}_k) - F(v_k)\bigr)$.
  \State Update SP-Broyden by $(d^{(r)}_k,\, g^{(r)}_k)$:
         $B_k^{(r+1)} \;=\; \mathsf{SP}\bigl(B_k^{(r)},\, d^{(r)}_k,\, g^{(r)}_k\bigr)$
         \quad ($B_k^{(0)} := B_k$).
  \State Compute next candidate
         $d^{(r+1)}_k \;=\; -(B_k^{(r+1)} + \beta_k I)^{-1} F(v_k)$.
  \If{$\Norm{d^{(r+1)}_k - d^{(r)}_k} \;\le\; \gamma\,\Norm{d^{(r+1)}_k}^{2}$}
       \State \textbf{break} \Comment{condition~\eqref{eq:cond14} verified, see Lemma~\ref{lem:app-afd:trigger}}
  \EndIf
\EndFor
\State Set $r^{*} := r$, $x_k := v_k + d^{(r^{*})}_k$.
\State \Return $x_k$, $B_{k+1} := B_k^{(r^{*}+1)}$.
\end{algorithmic}
\end{algorithm}

Существенные числовые константы~--- $\tau$, $\gamma$, $r_{\max}$~---
выбираются позже единообразно для всех $k, s$.

\section{Покомпонентные оценки}
\label{sec:app-afd:components}

В этом разделе устанавливаются две леммы, на которые опирается основное
утверждение, в дополнение к уже доказанной лемме~\ref{lem:fd}
главы~\ref{ch:afd} (FD-направленная производная,~---
стандартная оценка остатка формулы Ньютона--Лейбница~\cite[Lemma~4.1.12]{dennis1996}):
для любого $d \ne 0$ при $h = \tau \Norm{d}$, $\tau \in (0, 1]$, FD-приближение
$g := \tfrac{1}{\tau}(F(v_k + \tau d) - F(v_k))$ удовлетворяет
$\Norm{g - J_k d} \le \tfrac{L_1 \tau}{2}\Norm{d}^2$.

\begin{lemma}[Контроль ошибки матрицы вдоль текущего шага]
\label{lem:app-afd:dirbound}
Пусть $B_k^{(r+1)}$ получено SP-обновлением по паре $(d^{(r)}_k, g^{(r)}_k)$ из
алгоритма~\ref{alg:afd-inner}. Тогда после обновления
\begin{equation}
\Norm{\bigl(B_k^{(r+1)} - J_k\bigr)\, d^{(r)}_k}
\;\le\;
\frac{L_1\, \tau}{2}\,\Norm{d^{(r)}_k}^{2}.
\label{eq:app-afd:dirbound}
\end{equation}
\end{lemma}

\begin{proof}
Секант-сохраняющее SP-обновление по паре $(d, g)$ обеспечивает
$B_k^{(r+1)} d^{(r)}_k = g^{(r)}_k$. Применяя лемму~\ref{lem:fd}
с $s = d^{(r)}_k$:
\[
\Norm{(B_k^{(r+1)} - J_k)\, d^{(r)}_k}
= \Norm{g^{(r)}_k - J_k\, d^{(r)}_k}
\le \tfrac{L_1 \tau}{2}\,\Norm{d^{(r)}_k}^{2}. \qed
\]
\end{proof}

\begin{remark}[Сохранение равномерной оценки $M_0$]
\label{rem:app-afd:M0-preserved}
Frobenius-расстояние $\Fro{B_k^{(r+1)} - J_k}$ при FD-обновлениях
сохраняет равномерную оценку $M_0$. Действительно, по структуре
SP-обновления (точная пифагоровская декомпозиция, см.\ теорему~\ref{thm:sp-bd})
$\Fro{B_k^{(r+1)} - J_k}^{2} \le \Fro{B_k^{(r)} - J_k}^{2} +
\widetilde C_{\mathrm{SP}}\Norm{d^{(r)}_k}^{2}$. Используя $\Norm{d^{(r)}_k}\le D$
и не более $r_{\max}$ refinement-шагов на итерацию, итоговый прирост
ограничен $r_{\max}\widetilde C_{\mathrm{SP}}D^{2}$~--- постоянная величина,
поглощаемая в определении $M_0$ преамбулы (с учётом дополнительного
слагаемого $L_1 D \sqrt{d}$ от триангуляции $\Fro{\Jstar - J_k}$).
\end{remark}

\begin{lemma}[Триггер выхода из refinement-цикла]
\label{lem:app-afd:trigger}
Пусть $\Fro{B_k^{(r+1)} - J_k} \le M_0$ для всех refinement-шагов и пусть
триггер $\Norm{d^{(r+1)}_k - d^{(r)}_k} \le \gamma\,\Norm{d^{(r+1)}_k}^{2}$
из строки~6 алгоритма~\ref{alg:afd-inner} срабатывает. Обозначим
$M_J := M_0 + J^{*}_{\max} + L_1 D$~--- равномерную операторную оценку
$\Norm{B_k^{(r+1)} - J_k}_{op}$. Если
\begin{equation}
\gamma \;\le\; \frac{L_1}{4 M_J}
\quad\text{и}\quad
\tau \;\le\; \frac{1}{4 (1 + \gamma\Norm{d^{(r+1)}_k})^{2}},
\label{eq:app-afd:gamma-tau}
\end{equation}
то условие~\eqref{eq:cond14} выполнено для шага $d^{(r+1)}_k$:
\begin{equation}
\Norm{\bigl(B_k^{(r+1)} - J_k\bigr)\, d^{(r+1)}_k}
\;\le\;
\frac{L_1}{2}\,\Norm{d^{(r+1)}_k}^{2}.
\label{eq:app-afd:cond14}
\end{equation}
\end{lemma}

\begin{proof}
Раскладываем
\[
(B_k^{(r+1)} - J_k)\, d^{(r+1)}_k
\;=\;
(B_k^{(r+1)} - J_k)\, d^{(r)}_k
\;+\;
(B_k^{(r+1)} - J_k)\,(d^{(r+1)}_k - d^{(r)}_k).
\]
Первое слагаемое оценивается по лемме~\ref{lem:app-afd:dirbound}:
$\le (L_1\tau/2)\Norm{d^{(r)}_k}^{2}$.
Второе~--- через операторную норму
$\Norm{B_k^{(r+1)} - J_k}_{op} \le \Fro{B_k^{(r+1)} - \Jstar}
+ \Norm{\Jstar - J_k}_{op} \le \tilde M_{0} + L_1\Norm{v_k - x^{*}}
\le \tilde M_{0} + L_1 D \le M_J$,
затем по триггеру:
$\le M_J\, \gamma\, \Norm{d^{(r+1)}_k}^{2}$.

Используя $\Norm{d^{(r)}_k} \le \Norm{d^{(r+1)}_k} + \Norm{d^{(r+1)}_k - d^{(r)}_k}
\le (1 + \gamma \Norm{d^{(r+1)}_k})\,\Norm{d^{(r+1)}_k}$, имеем
$\Norm{d^{(r)}_k}^{2} \le (1 + \gamma\Norm{d^{(r+1)}_k})^{2}\, \Norm{d^{(r+1)}_k}^{2}$.
Складывая два слагаемых:
\[
\Norm{(B_k^{(r+1)} - J_k)\, d^{(r+1)}_k}
\;\le\;
\bigl[\tfrac{L_1\tau}{2}(1 + \gamma\Norm{d^{(r+1)}_k})^{2}
        + M_J\gamma\bigr]\,\Norm{d^{(r+1)}_k}^{2}.
\]
По~\eqref{eq:app-afd:gamma-tau} коэффициент в скобках $\le L_1/4 + L_1/4 = L_1/2$,
что даёт~\eqref{eq:app-afd:cond14}. \qed
\end{proof}

\begin{remark}[Калибровка $\tau, \gamma$]
\label{rem:app-afd:calib}
В терминальной фазе $\Norm{d^{(r+1)}_k} \to 0$, поэтому условие
$\gamma\Norm{d^{(r+1)}_k} \le 1$ из~\eqref{eq:app-afd:gamma-tau} при
любом фиксированном $\gamma$ (не зависящем от $k$) выполнено
автоматически, и $(1 + \gamma\Norm{\cdot})^{2} \le 4$. Значение по умолчанию
$\tau = 1/8$ при этом обеспечивает $\tau\cdot 4 = 1/2$, т.\,е.\
$L_1\tau/2 \cdot 4 = L_1/4$, что в сумме с $M_J\gamma = L_1/4$ даёт ровно
$L_1/2$.
\end{remark}

\section{Конечность refinement-цикла}
\label{sec:app-afd:rstar}

Покажем, что число refinement-шагов $r^{*}$ ограничено сверху и в
терминальной фазе сходится к $r^{*} = 1$ (одна FD-секущая в среднем).

\begin{proposition}[Ограниченность $r^{*}$]
\label{prop:app-afd:rstar}
Алгоритм~\ref{alg:afd-inner} с гиперпараметром $r_{\max} \in \NN$
удовлетворяет $r^{*}(k, s) \le r_{\max}$ для всех $k, s$ по построению
(трункация цикла~$\mathbf{for}$). При срабатывании триггера до $r_{\max}$
последняя матрица $B_k^{(r^{*}+1)}$ удовлетворяет~\eqref{eq:cond14}
по лемме~\ref{lem:app-afd:trigger}; при $r^{*} = r_{\max}$
условие~\eqref{eq:cond14} в общем случае не гарантировано, и шаг
обрабатывается через rate~\cite[Theorem~3.2]{Agafonov2024VIJI}
($\mathcal O(T^{-1})$ вместо $\mathcal O(T^{-3/2})$); это лишь ужесточает
префактор глобальной асимптотики на постоянный множитель.
\end{proposition}

\begin{remark}[Эмпирическое поведение $r^{*}$]
\label{rem:app-afd:rstar-emp}
На гладких сильно монотонных задачах эмпирический $r^{*}$ часто выходит
на 1 в окрестности корня (рис.~\ref{fig:viji_rstar} основного текста);
структурная причина~--- alignment SP-Broyden-обновления вдоль текущего
шага. Это наблюдение не является теоретической гарантией в общей форме:
в неасимптотическом режиме $\Norm{d_k^{(r+1)} - d_k^{(r)}}$ масштабируется
линейно по $\Norm{e_k}$ (через постоянный член $M_0/\sigma_0$ в
SP-обновлении), а правая часть триггера~--- квадратично, поэтому
гарантированное срабатывание триггера обеспечивается лишь алгоритмической
трункацией $r^{*} \le r_{\max}$.
\end{remark}

\section{Сборка глобальной скорости теоремы~\ref{thm:afd}}
\label{sec:app-afd:cond14}

\begin{proof}[Доказательство теоремы~\ref{thm:afd}]
Зафиксируем $\tau = 1/8$ и $\gamma = L_1/(4 M_J)$ с
$M_J = M_0 + J^{*}_{\max} + L_1 D$ согласно
лемме~\ref{lem:app-afd:trigger}. Тогда оба
условия~\eqref{eq:app-afd:gamma-tau} выполнены для всех $k, s$ при
$\gamma\Norm{d^{(r+1)}_k} \le 1$ (что эквивалентно $\Norm{d^{(r+1)}_k}
\le 4 M_J/L_1$ и автоматически выполнено в терминальной фазе).

\textbf{Итерации со сработавшим триггером} ($r^{*} \le r_{\max}-1$).
По лемме~\ref{lem:app-afd:trigger} матрица $B_k^{(r^{*}+1)}$
удовлетворяет~\eqref{eq:cond14} с константой $L_1/2$. Применяя
theorem~3.3 из~\cite{Agafonov2024VIJI} к каждой такой внутренней итерации
сегмента, получаем GAP $= \mathcal O(L_1 D_s^{3} / N^{3/2})$.

\textbf{Итерации с трункацией} ($r^{*} = r_{\max}$).
По proposition~\ref{prop:app-afd:rstar} в общем случае условие~\eqref{eq:cond14}
не гарантировано; такие шаги обрабатываются через theorem~3.2
из~\cite{Agafonov2024VIJI} с rate $\mathcal O(L_1 D_s^{2}/N)$. На асимптотику
сегмента это влияет лишь как постоянный коэффициент в префакторе.

\textbf{Глобальная сходимость.}
Сегментная геометрическая декада~\eqref{eq:restart-geom} с фактором
$\rho \in (0, 1)$ даёт $D_S \le \rho^S D_0$, и суммарно
$\mathrm{GAP}(\bar x^{(S)}_N) = \mathcal O(L_1 D_0^{3} \rho^{3S} / N^{3/2})$.
Стоимость одной внутренней итерации AFD~--- $1 + r_{\max}$ вызовов $F$
(основной шаг $+$ до $r_{\max}$ refinement-шагов). Переход к $\#F$:
\[
\#F\bigl(\mathrm{GAP} \le \varepsilon\bigr)
\;=\;
(1+r_{\max})\cdot\mathcal O\!\bigl((L_1 D_0^{3}/\varepsilon)^{2/3}\bigr),
\]
что и есть утверждение теоремы~\ref{thm:afd}. \qed
\end{proof}

\section{Стоимость и сравнение с VIQA-Broyden}
\label{sec:app-afd:cost}

Стоимость одной внутренней итерации AFD~--- $1 + r_{\max}$ вызовов $F$
(постоянный гиперпараметр алгоритма). По сравнению с одним вызовом $F$
у VIQA-Broyden это даёт постоянный множитель $1+r_{\max}$ в $F$-сложности,
но при \emph{оптимальной} скорости $\mathcal O(T^{-3/2})$ вместо
неоптимальной $\mathcal O(T^{-1})$. Конкретно, для достижения точности
$\varepsilon$ в GAP:
\[
\#F_{\mathrm{AFD}}(\varepsilon)
\;=\;
(1+r_{\max})\cdot \mathcal O\bigl((L_1 D^{3}/\varepsilon)^{2/3}\bigr),
\qquad
\#F_{\mathrm{VIQA\text{-}Broyden}}(\varepsilon)
\;=\;
\mathcal O(L_1 D^{3}/\varepsilon),
\]
что даёт асимптотическое преимущество AFD $\sim (L_1 D^{3}/\varepsilon)^{1/3}$
раз~--- неограниченно растущее с уменьшением $\varepsilon$.

\begin{remark}[Чувствительность к $\tau$]
\label{rem:app-afd:tau-sens}
Лемма~\ref{lem:app-afd:trigger} показывает, что $\tau$ может быть
выбрано в диапазоне $\tau \in (0, 1/8]$ без изменения асимптотики.
Уменьшение $\tau$ снижает первый член $(L_1\tau/2)\Norm{d^{(0)}_k}^{2}$
за счёт пропорционального уменьшения $\Norm{F(v_k + \tau d) - F(v_k)}$,
что повышает чувствительность к машинной точности при
$\Norm{d_k} \lesssim \sqrt{\varepsilon_{\mathrm{mach}}}$. Значение по
умолчанию $\tau = 1/8$ обеспечивает в-точности константу $L_1/2$
в~\eqref{eq:app-afd:cond14} с балансом <<систематика vs.\ шум>>.
\end{remark}

\paragraph{Сводное замечание.}
Доказанная теорема показывает, что AFD-обновление якобиана с трункацией
refinement-цикла на $r_{\max}$ обеспечивает оптимальную скорость
$\mathcal O(T^{-3/2})$ для VIJI-Restarted на $\mu$-сильно монотонных
$L_1$-гладких VI без аналитического якобиана~--- ценой постоянного
коэффициента $1+r_{\max}$ в $F$-сложности.
